% notes-sequent-calculus-lk.tex
% Topic: Sequent Calculus LK

\begin{remark*}[Exposition]
This section fixes Gentzen's classical first-order sequent calculus LK in the
style of Takeuti's \emph{Proof Theory}. The goal is not yet to prove the
Hauptsatz, but to establish the formal arena in which proof transformations
can be stated without ambiguity.
\end{remark*}

\begin{remark*}[Notation]
Two arrows are kept distinct throughout this section. The symbol $\supset$ is
the object-language implication connective inside formulas. The symbol $\to$
is the metalanguage sequent arrow separating antecedent from succedent. Thus
$A \supset B$ is a formula, while $\Gamma \to \Delta$ is a sequent.
\end{remark*}

\subsection{Language}

\begin{definitionbox}{Definition (LK Alphabet)}
\begin{definition}[LK Alphabet]\label{def:lk-alphabet}
An \textbf{LK alphabet} consists of the following disjoint syntactic classes:
\begin{enumerate}
  \item a denumerable set $FV = \{a_0,a_1,a_2,\ldots\}$ of free variables;
  \item a denumerable set $BV = \{x_0,x_1,x_2,\ldots\}$ of bound variables,
        disjoint from $FV$;
  \item for each $n \geq 0$, a denumerable set of $n$-ary function symbols;
  \item for each $n \geq 0$, a denumerable set of $n$-ary predicate symbols;
  \item logical connectives $\lnot,\land,\lor,\supset$;
  \item quantifier symbols $\forall,\exists$;
  \item punctuation symbols for grouping and lists.
\end{enumerate}
\end{definition}
\end{definitionbox}

\begin{remark*}[Standard quantified statement]
$FV$ and $BV$ are denumerable disjoint sets, each non-logical symbol has a
fixed arity, and the logical alphabet contains exactly the displayed
connectives, quantifiers, and punctuation symbols.
\end{remark*}

\begin{remark*}[Interpretation]
The split between free and bound variables is deliberate. Terms are built from
free variables, while quantifiers bind bound variables. This convention makes
the eigenvariable conditions in the quantifier rules visible at the level of
syntax rather than hiding them inside an informal substitution convention.
\end{remark*}

\DefinitionalRoot

\begin{definitionbox}{Definition (LK Term)}
\begin{definition}[LK Term]\label{def:lk-term}
The set $\mathrm{Term}_{LK}$ of \textbf{LK terms} is the smallest set
such that:
\begin{enumerate}
  \item if $a \in FV$, then $a \in \mathrm{Term}_{LK}$;
  \item if $f$ is an $n$-ary function symbol and
        $t_1,\ldots,t_n \in \mathrm{Term}_{LK}$, then
        $f(t_1,\ldots,t_n) \in \mathrm{Term}_{LK}$.
\end{enumerate}
In particular, a $0$-ary function symbol is a term.
\end{definition}
\end{definitionbox}

\begin{remark*}[Standard quantified statement]
$\mathrm{Term}_{LK}$ is the least set containing $FV$ and closed under
$n$-ary function-symbol application.
\end{remark*}

\begin{remark*}[Interpretation]
No bound variable is a term. This is one of Takeuti's bookkeeping choices:
substitution acts on free variables, and quantification is described by
promoting designated free-variable occurrences to fresh bound-variable
occurrences.
\end{remark*}

\begin{dependencies}
\begin{itemize}
  \item \hyperref[def:lk-alphabet]{LK Alphabet}
\end{itemize}
\end{dependencies}

\begin{definitionbox}{Definition (LK Formula)}
\begin{definition}[LK Formula]\label{def:lk-formula}
The set $\mathrm{Form}_{LK}$ of \textbf{LK formulas} is the smallest set
such that:
\begin{enumerate}
  \item if $R$ is an $n$-ary predicate symbol and
        $t_1,\ldots,t_n \in \mathrm{Term}_{LK}$, then
        $R(t_1,\ldots,t_n) \in \mathrm{Form}_{LK}$;
  \item if $A,B \in \mathrm{Form}_{LK}$, then
        $\lnot A$, $(A \land B)$, $(A \lor B)$, and
        $(A \supset B)$ belong to $\mathrm{Form}_{LK}$;
  \item if $A \in \mathrm{Form}_{LK}$, $a \in FV$, and
        $x \in BV$ is fresh for $A$, then
        $\forall x\,A[x/a]$ and $\exists x\,A[x/a]$ belong to
        $\mathrm{Form}_{LK}$.
\end{enumerate}
\end{definition}
\end{definitionbox}

\begin{remark*}[Standard quantified statement]
$\mathrm{Form}_{LK}$ is the least set containing the atomic formulas over
$\mathrm{Term}_{LK}$ and closed under $\lnot$, $\land$, $\lor$, $\supset$,
and the fresh-bound-variable quantifier clauses.
\end{remark*}

\begin{remark*}[Interpretation]
The formula clauses define the object language. Quantified formulas are built
from already formed formulas by replacing selected free-variable occurrences
with a fresh bound variable.
\end{remark*}

\begin{remark*}[Exposition]
When a formula is written $F(a)$, the displayed $a$ marks designated
occurrences of a free variable. For a term $t$, the notation $F(t)$ denotes
the result of replacing exactly those occurrences by $t$. If $x \in BV$ is
fresh, then $F(x)$ denotes the corresponding replacement by $x$, so that
$\forall x\,F(x)$ and $\exists x\,F(x)$ are formulas.
\end{remark*}

\begin{dependencies}
\begin{itemize}
  \item \hyperref[def:lk-alphabet]{LK Alphabet}
  \item \hyperref[def:lk-term]{LK Term}
\end{itemize}
\end{dependencies}

\subsection{Sequents and Interpretation}

\begin{definitionbox}{Definition (LK Sequent)}
\begin{definition}[LK Sequent]\label{def:lk-sequent}
An \textbf{LK sequent} is an expression
\[
  \Gamma \to \Delta,
\]
where $\Gamma = A_1,\ldots,A_m$ and $\Delta = B_1,\ldots,B_n$ are finite,
possibly empty, sequences of LK formulas. The sequence $\Gamma$ is the
\textbf{antecedent}, and the sequence $\Delta$ is the \textbf{succedent}.
\end{definition}
\end{definitionbox}

\begin{remark*}[Standard quantified statement]
A sequent is an element of $\mathrm{Form}_{LK}^{*}\times
\mathrm{Form}_{LK}^{*}$, written $\Gamma \to \Delta$.
\end{remark*}

\begin{remark*}[Interpretation]
A sequent is not itself a formula. It is a pair of finite formula lists. The
intended reading is that the conjunction of the antecedent formulas implies
the disjunction of the succedent formulas. Repetitions and order are retained
at the syntactic level because the structural rules explicitly manipulate
weakening, contraction, and exchange.
\end{remark*}

\begin{remark*}[Exposition]
The formula associated to a sequent is
\[
\iota(\Gamma \to \Delta) =
\begin{cases}
(A_1 \land \cdots \land A_m) \supset (B_1 \lor \cdots \lor B_n),
  & m \geq 1,\ n \geq 1,\\
B_1 \lor \cdots \lor B_n,
  & m = 0,\ n \geq 1,\\
\lnot(A_1 \land \cdots \land A_m),
  & m \geq 1,\ n = 0,\\
\bot,
  & m = 0,\ n = 0.
\end{cases}
\]
The sequent is valid exactly when this associated formula is valid in the
usual first-order semantic sense.
\end{remark*}

\begin{dependencies}
\begin{itemize}
  \item \hyperref[def:lk-formula]{LK Formula}
\end{itemize}
\end{dependencies}

\subsection{Inferences and Rules}

\begin{definitionbox}{Definition (LK Inference)}
\begin{definition}[LK Inference]\label{def:lk-inference}
An \textbf{LK inference} is an instance of one of the LK rule schemata with
one or two upper sequents and one lower sequent. In a logical inference, the
\textbf{principal formula} is the formula introduced in the lower sequent, the
\textbf{auxiliary formula} or formulas are its immediate upper-sequent
components, and the remaining formulas are \textbf{side formulas}.
\end{definition}
\end{definitionbox}

\begin{remark*}[Standard quantified statement]
An inference is a rule instance whose lower sequent follows from one or two
upper sequents according to one of the LK structural, cut, connective, or
quantifier schemata.
\end{remark*}

\begin{remark*}[Interpretation]
Inference rules describe how proof trees grow downward. The principal formula
records what has just been introduced; side formulas record the context
carried unchanged through the rule.
\end{remark*}

\begin{remark*}[Exposition]
LK has left and right weakening, contraction, and exchange:
\[
\frac{\Gamma \to \Delta}{A,\Gamma \to \Delta}\ (w:l)
\qquad
\frac{\Gamma \to \Delta}{\Gamma \to \Delta,A}\ (w:r)
\]
\[
\frac{A,A,\Gamma \to \Delta}{A,\Gamma \to \Delta}\ (c:l)
\qquad
\frac{\Gamma \to \Delta,A,A}{\Gamma \to \Delta,A}\ (c:r)
\]
\[
\frac{\Gamma,A,B,\Pi \to \Delta}{\Gamma,B,A,\Pi \to \Delta}\ (e:l)
\qquad
\frac{\Gamma \to \Delta,A,B,\Lambda}{\Gamma \to \Delta,B,A,\Lambda}\ (e:r).
\]
These are rules because sequents are sequences. A multiset presentation would
hide exchange, and a set presentation would also hide contraction and some
weakening behavior.
\end{remark*}

\begin{remark*}[Exposition]
The cut rule is
\[
\frac{\Gamma \to \Delta,A \qquad A,\Pi \to \Lambda}
     {\Gamma,\Pi \to \Delta,\Lambda}\ (cut).
\]
The formula $A$ is the cut formula. Takeuti treats cut as structural, but it is
distinguished because Gentzen's Hauptsatz says that cuts can be eliminated
from LK proofs.
\end{remark*}

\begin{remark*}[Exposition]
Each connective and quantifier has a left rule, introducing the connective or
quantifier in the antecedent, and a right rule, introducing it in the
succedent. For example:
\[
\frac{\Gamma \to \Delta,A}{\lnot A,\Gamma \to \Delta}\ (\lnot:l)
\qquad
\frac{A,\Gamma \to \Delta}{\Gamma \to \Delta,\lnot A}\ (\lnot:r),
\]
\[
\frac{A,\Gamma \to \Delta}{A \land B,\Gamma \to \Delta}\ (\land:l_1)
\qquad
\frac{B,\Gamma \to \Delta}{A \land B,\Gamma \to \Delta}\ (\land:l_2)
\qquad
\frac{\Gamma \to \Delta,A \qquad \Gamma \to \Delta,B}
     {\Gamma \to \Delta,A \land B}\ (\land:r),
\]
\[
\frac{A,\Gamma \to \Delta \qquad B,\Gamma \to \Delta}
     {A \lor B,\Gamma \to \Delta}\ (\lor:l)
\qquad
\frac{\Gamma \to \Delta,A}{\Gamma \to \Delta,A \lor B}\ (\lor:r_1)
\qquad
\frac{\Gamma \to \Delta,B}{\Gamma \to \Delta,A \lor B}\ (\lor:r_2),
\]
\[
\frac{\Gamma \to \Delta,A \qquad B,\Pi \to \Lambda}
     {A \supset B,\Gamma,\Pi \to \Delta,\Lambda}\ (\supset:l)
\qquad
\frac{A,\Gamma \to \Delta,B}{\Gamma \to \Delta,A \supset B}\ (\supset:r).
\]
\end{remark*}

\begin{remark*}[Exposition]
For quantifiers, $F(t)$ denotes substitution into designated occurrences:
\[
\frac{F(t),\Gamma \to \Delta}
     {\forall x\,F(x),\Gamma \to \Delta}\ (\forall:l)
\qquad
\frac{\Gamma \to \Delta,F(a)}
     {\Gamma \to \Delta,\forall x\,F(x)}\ (\forall:r),
\]
\[
\frac{F(a),\Gamma \to \Delta}
     {\exists x\,F(x),\Gamma \to \Delta}\ (\exists:l)
\qquad
\frac{\Gamma \to \Delta,F(t)}
     {\Gamma \to \Delta,\exists x\,F(x)}\ (\exists:r).
\]
In $(\forall:r)$ and $(\exists:l)$, the eigenvariable $a$ must not occur in
the lower sequent. This condition expresses that $a$ is arbitrary in
$(\forall:r)$ and a fresh name for a witness in $(\exists:l)$.
\end{remark*}

\begin{dependencies}
\begin{itemize}
  \item \hyperref[def:lk-formula]{LK Formula}
  \item \hyperref[def:lk-sequent]{LK Sequent}
\end{itemize}
\end{dependencies}

\subsection{Proofs}

\begin{definitionbox}{Definition (LK-Proof)}
\begin{definition}[LK-Proof]\label{def:lk-proof}
An \textbf{LK-proof} is a finite tree of sequents such that every leaf is an
initial sequent of the form $D \to D$, and every internal node is obtained
from its child node or child nodes by an LK inference. The label of the root is
the \textbf{endsequent}. A sequent $S$ is \textbf{LK-provable}, written
$\vdash_{LK} S$, when $S$ is the endsequent of some LK-proof.
\end{definition}
\end{definitionbox}

\begin{remark*}[Standard quantified statement]
$P$ is an LK-proof of $S$ iff $P$ is a finite sequent-labeled tree whose
leaves are initial sequents, whose internal nodes are LK inferences, and whose
root is labeled $S$.
\end{remark*}

\begin{remark*}[Interpretation]
An LK-proof is a certified derivation of its endsequent from initial identity
sequents. The rule labels explain how each lower sequent is obtained from its
upper sequent or sequents.
\end{remark*}

\begin{remark*}[Exposition]
Drawn Gentzen-style, proof trees have their initial sequents at the top and
their endsequent at the bottom. A branch is a path from the endsequent upward
to an initial sequent. A thread follows a formula occurrence through the tree;
threads become important when organizing cut-elimination arguments.
\end{remark*}

\begin{dependencies}
\begin{itemize}
  \item \hyperref[def:lk-sequent]{LK Sequent}
  \item \hyperref[def:lk-inference]{LK Inference}
\end{itemize}
\end{dependencies}

\begin{definitionbox}{Definition (LJ-Proof)}
\begin{definition}[LJ-Proof]\label{def:lj-proof}
An \textbf{LJ-proof} is an LK-proof in which every sequent has succedent length
at most one.
\end{definition}
\end{definitionbox}

\begin{remark*}[Standard quantified statement]
$P$ is an LJ-proof iff $P$ is an LK-proof and every sequent $\Gamma \to
\Delta$ occurring in $P$ satisfies $|\Delta| \leq 1$.
\end{remark*}

\begin{remark*}[Interpretation]
LK is classical first-order logic. LJ is Gentzen's intuitionistic sequent
calculus obtained by restricting every succedent to length at most one. This
single structural restriction removes the general classical use of multiple
alternative conclusions.
\end{remark*}

\begin{remark*}[Exposition]
Gentzen's Hauptsatz states that every LK-proof, and similarly every LJ-proof,
can be transformed into a proof of the same endsequent without using cut. The
subformula property, consistency results, interpolation phenomena, and
Herbrand-style consequences all begin from this cut-elimination theorem.
\end{remark*}

\begin{dependencies}
\begin{itemize}
  \item \hyperref[def:lk-sequent]{LK Sequent}
  \item \hyperref[def:lk-proof]{LK-Proof}
\end{itemize}
\end{dependencies}
